% !TEX program = xelatex
\documentclass[UTF8,fontset=fandol,zihao=-4,a4paper]{ctexart}
\usepackage[margin=22mm]{geometry}
\usepackage{amsmath,amssymb,amsthm,mathtools,booktabs,array,longtable}
\usepackage{graphicx,xcolor,listings,hyperref}
\hypersetup{colorlinks=true,linkcolor=blue!50!black,urlcolor=blue!60!black,
 pdftitle={五维中心方体极大算子的一个准确弱型上界},
 pdfauthor={五维中心方体研究项目}}
\setlength{\emergencystretch}{2em}
\newtheorem{theorem}{定理}[section]
\newtheorem{lemma}[theorem]{引理}
\newtheorem{proposition}[theorem]{命题}
\newtheorem{remark}[theorem]{注}
\newcommand{\R}{\mathbb R}
\newcommand{\E}{\mathbb E}
\newcommand{\1}{\mathbf 1}
\newcommand{\dd}{\,\mathrm d}
\newcommand{\norm}[1]{\left\lVert#1\right\rVert}
\DeclareMathOperator{\supp}{supp}
\DeclareMathOperator{\essinf}{ess\,inf}
\lstset{basicstyle=\ttfamily\footnotesize,breaklines=true,columns=fullflexible,
 frame=single,keepspaces=true}
\title{五维中心方体极大算子的一个准确弱型上界\\
\large 有限正源、正时间曲率与可复核积分证书}
\author{五维中心方体研究项目\\\small 内部研究稿}
\date{2026年10月27日}
\begin{document}
\maketitle
\begin{abstract}
研究固定Lebesgue测度下五维中心轴平行方体极大算子的弱$(1,1)$常数。
所有坐标共用一个任意正半径，输入为一般$L^1(\R^5)$函数。
以十六方向Cauchy半群的Green势为比较核，证明全空间有限正源及同一接收集的全尺度传递。
通过完整正面覆盖、固定锚点曲率、逆时间正矩与准确凸弦支付，获得
\[
 C_5\le U_{\rm cert}\approx28.936078340172045<28.936079.
\]
展示小数不参与认证。论文给出短有理常数与一个纯整数充分条件；完整证书使用原始有理数。
五维奇异直线、平面及其交叠通过全空间分布论和完整小时间预算处理。
两轮审计采用不同的整数及有理运算路径，分别核验完整积分与完整空间覆盖，
并明确原始$H_\epsilon$和物理导数区间仍是继承的严格数值输入。
本稿的贡献为上界；既有内部下界及尚未闭合的局部候选单独标识。
\end{abstract}

\section{对象、比较门槛与主结果}
令$Q=[-1,1]^5$，定义
\[
 \mathcal M_5f(x)=\sup_{r>0}\frac1{(2r)^5}
      \int_{x+[-r,r]^5}|f(y)|\dd y,\qquad
 C_5=\sup_{\substack{f\in L^1,\ f\ne0\\\lambda>0}}
 \frac{\lambda|\{\mathcal M_5f>\lambda\}|}{\norm f_1}.
\]
这里的$r$对五个坐标相同。以下所有上界均支付一般输入及全部正尺度。

Aldaz的固定维构造给出五维显式公开下界$A_5\approx1.7481246343$；
对更大背景测度类的统一方体估计给出$C_5\le32$；
维数单调性允许把已有四维下界传递至五维\cite{Aldaz2000,Aldaz2011,Aldaz2020}。
本项目在2026年10月25日至26日实际核查的原始来源中，采用的公开比较区间为
$[A_5,32]$。这是一项有范围的文献核查，不是对全球文献的穷尽性声明。
部分扫掠的比较思想亦见作者原始方法网页\cite{Lin2026}；
该网页列出的平面方体结论不能直接代入五维。

\begin{theorem}\label{thm:main}
存在本稿所述完整准确证书，使
\[
 C_5\le U_{\rm cert}<\frac{28936079}{10^6}=28.936079.
\]
更准确地，令$g_*$、$H_*$、$J_*$为附录中的正有理数，则
\[
 U_{\rm cert}=\left(\frac{3125(H_*+J_*)^4}{g_*^5}\right)^{1/4},
\]
取唯一正四次根。其准确有理隔离及新旧差值见随文证书。
\end{theorem}

\begin{table}[ht]
\centering\small
\begin{tabular}{p{.29\textwidth}p{.27\textwidth}p{.33\textwidth}}
\toprule
比较对象&下界&上界\\ \midrule
已核查的公开来源&$A_5$&$32$\\
原开工内部门槛&继承四维完整证书&$32$\\
冻结接续门槛&$C_{\rm cert}$&$32$\\
本稿完整结果&继承$C_{\rm cert}$&$U_{\rm cert}<28.936079$\\
\bottomrule
\end{tabular}
\caption{完整门槛。$C_{\rm cert}\approx1.9943164615684859$为既有内部证书，
其准确代数定义与审计路径保留在账本中。本轮没有提高接续下界。}
\end{table}

本稿不将内部结果称为公开发表成果。下界研究已按用户指令停止。
接续门槛要求的双侧突破尚未成立，双侧突破后的三轮迭代也未触发。
根据后续明确要求，本稿交付已经完成双轮审计的当前上界成果。

\section{辅助半群与一般输入的有限源传递}
\subsection{归一化和势核}
取
\[
 V=\{(1,\pm1,\pm1,\pm1,\pm1)\},\qquad
 \psi(\xi)=\sum_{v\in V}|v\cdot\xi|,\qquad
 \widehat p_t(\xi)=e^{-t\psi(\xi)}.
\]
方向不单位化，各项权为$1$、时间均为$t$；逆Fourier因子为$(2\pi)^{-5}$。
若$B$为这十六列组成的矩阵，则$BB^{T}=16I_5$，
从而$4|\xi|\le\psi(\xi)\le16|\xi|$。
全$1$方向与尾部恰有一个$+1$的四方向组成的五方向基，行列式绝对值为$48$，故
\[
 0<p_t(x)\le\frac1{48\pi^5t^5},\qquad \int_{\R^5}p_t=1.
\]
设$P_t f=p_t*f$，采用\emph{负的Markov生成元}约定
$\widehat{Lf}=-\psi\widehat f$；其$L^2$定义域为$H^1$。
记
\[
 G=\int_0^\infty p_t\dd t,\quad
 H_\epsilon=\int_\epsilon^\infty p_t\dd t,\quad
 F=\{(1,y):y\in[0,1]^4\},\quad p=\frac54.
\]
完整$G$为$-4$齐次扩展值正函数；固定截时$H_\epsilon$不具有该齐次性。

热核的质量与峰值拆分给$G\in L^{5/4,\infty}(\R^5)$：
对$T>0$，$\norm{\int_0^Tp_t\dd t}_1=T$，
而$\norm{\int_T^\infty p_t\dd t}_\infty\le C T^{-4}$。
据此对任意$c>0$，$W=(G-c)_+\in L^1$。

\begin{lemma}[全空间有限正源]\label{lem:source}
分布$\mu=\delta_0+LW$是正Radon测度，
$\mu(\R^5)=1$且$\mu(\{0\})=0$。
\end{lemma}
\begin{proof}
设$U_t=\int_0^tp_s\dd s$。Tonelli给$P_tG=\int_t^\infty p_s\dd s$；
在$G$有限的满测集上可写$P_tG=G-U_t$。由正部凸性，
$P_tW\ge(G-U_t-c)_+\ge W-U_t$。
用全空间非负测试函数取$t\downarrow0$，得到$LW+\delta_0\ge0$。
此论证使用全空间测试函数；不在两个势同时无穷处作相减，也没有删除奇异子空间或水平集边界。
光滑大截断$\chi_R$满足$\norm{L\chi_R}_\infty=O(R^{-1})$，
所以$\int\chi_R\dd\mu\to1$；正性给总质量$1$。
又$LG=-\delta_0$，因而$\mu=-L\min(G,c)$。
原点附近的小截断满足$\norm{L\chi_\eta}_1=O(\eta^4)$，
而$\min(G,c)$有界，故$\mu(\{0\})=0$。
这些生成元估计由每个一维跳跃方向的二差分公式分别在小、大参数区间得到。
\end{proof}

\begin{proposition}[有限源比较]\label{prop:transfer}
若$g>0$、$G\ge g$于$Q$几乎处处，且
$I=\int_{[0,1]^4}G(1,y)^{5/4}\dd y<\infty$，则
\[
 C_5\le5^{5/4}I\,g^{-5/4}.
\]
\end{proposition}
\begin{proof}
先取非负、有界、紧支撑$f$，质量$m$。对$\kappa>0$构造障碍势$u$，
使$u\ge0$、$\sigma=f+Lu$、$0\le\sigma\le\kappa$，
且$\sigma=\kappa$于$D=\{u>0\}$，$\int\sigma=m$。
下面给出确保其在本模型合法的构造要点。

固定$t>0$，从零迭代
$u_{k+1}=(P_tu_k+t(f-\kappa))_+$，单调极限为$u_t$。
以$V_t=t\sum_{j\ge0}P_{jt}f$控制，得
$u_t\le(V_t-\kappa t)_+$。
$V_t-tf\in C_0$使固定$t$的$u_t$可积；
热核峰值给统一上界$u_t\le A<\infty$。
令$L_t=(P_t-I)/t$，则$\sigma_t=f+L_tu_t$满足上述离散互补条件和质量等式。
关键统一估计为
\[
 \kappa\int u_t=\int fu_t+\langle u_t,L_tu_t\rangle\le mA.
\]
故没有把通常不可积的$V_t$误作统一$L^1$上界。
$\norm{L_tu_t}_2$有界；频率截断后的局部Rellich紧性及高频尾控制
给局部强$L^2$和几乎处处子列。取极限得$u\in L^1\cap H^1$及$\sigma$。
互补条件由局部收敛保留；大截断恢复$\int\sigma=m$，于是$\kappa|D|\le m$。

取$0<c<g$，置$N=32(g-c)$、$K=W/N$、
$K_r(x)=r^{-5}K(x/r)$。则
$K_r\ge(2r)^{-5}\1_{[-r,r]^5}$。
由引理\ref{lem:source}及Fourier恒等式，
\[
 K_r*Lu=\frac{\mu_r*u-u}{Nr},
\]
其中$\mu_r$为$\mu$的尺度推前。等式先在光滑函数成立，
再由$H^1$与有限测度的$L^2$卷积界延拓。
同一个$u,\sigma,D$对所有$r$适用。在$D$外$u=0$、$\mu_r*u\ge0$，故
$K_r*f\le K_r*\sigma\le\kappa\norm K_1$。
先合并有理$r>0$的可数零集；$f$有界时，
$\norm{K_r*f-K_s*f}_\infty\le\norm f_\infty\norm{K_r-K_s}_1$
给全部正$r$。因此弱型常数不超过$\norm K_1$。
对一般$f$用$\min(|f|,j)\1_{[-j,j]^5}\uparrow|f|$及单调收敛即可。

最后，面分解和完整$G$的齐次性给
$\norm W_1=128I c^{-1/4}$，故
$\norm K_1=4I c^{-1/4}/(g-c)$。
在$c=g/5$处最小，得到结论。这里的半径传递和归一化均未改变。
\end{proof}

\section{五维奇异区域与完整小时间预算}
\subsection{不依赖四维线外连续性的处理}
设$S_1$为十六条方向直线之并，$S_2$为方向对张成平面之并。
$F\cap S_1$为全$1$顶点，而
\[
 F\cap S_2=\bigcup_{\varnothing\ne J\subset\{2,3,4,5\}}
 \{y_j=s\ (j\in J),\ y_j=1\ (j\notin J),\ 0\le s\le1\}.
\]
这是十五条线段。写$x=av+bw$，首坐标给$a+b=1$；
同号坐标由正面条件固定为$1$，异号坐标等于$\pm(a-b)$。
混合符号只能在$s=0$处进入正面，得到该分类。
迹分类本身不是到全空间平面的距离公式。

对于任意投影子空间$E$，正推前与
$z^2t/[\pi(t^2+z^2)]\le t/\pi$给
\[
 |x|^2q_t^E(x)\le\frac{16t}{\pi}
 \sum_{P_Ev\ne0}|P_Ev|^{-1}q_t^{E\cap v^\perp}(P_{E\cap v^\perp}x).
 \tag{3.1}
\]
先用非负测试与Tonelli证明，随后以正时间密度连续性取得处处版本；
没有调用发散的Cauchy二阶矩。

本模型的准确多面体证书给出$\{\psi\le1\}$的全部280个真顶点。
其平方范数分布依次为$1/128$共40个、$1/144$共80个、$1/162$共160个。
平方范数的凸性把顶点上界延至整个凸包，故$\psi(\xi)\ge8\sqrt2|\xi|$。
完整顶点及扇覆盖的资格均保留在原准确资料中。
Fourier径向积分于是给
\[
 P_4=\frac3{65536\pi^2},\qquad P_3=\frac1{1024\sqrt2\pi^2}
\]
作为四维、三维投影峰值系数。

令$F_5=16/(\pi\sqrt5)$、
$S=5\sqrt5/4+5\sqrt{5/6}$、
$C_3=F_5P_4$、$C_1=480F_5SP_3/\pi$。
非中央十五方向在整个闭$F$上的距离平方至少$1/2$，
由(3.1)给总贡献$C_1t^{-1}$。
中央方向$v_+=(1,1,1,1,1)$按
$r=|P_{v_+^\perp}x|$分区。$r\le\eta=1/10$时
$|x|^2\ge5-9\eta+5\eta^2=83/20$；
$r\ge\eta$时中央贡献至多$C_1t^{-1}/(30\eta^2)$。
准确常数门
\[
 \frac{C_1\epsilon^2}{30\eta^2}\le\frac{C_3}{83/20},
 \qquad \epsilon=\frac1{512}
\]
已严格认证，故全$F$、全部$0<t\le\epsilon$有
\[
 \begin{gathered}
 p_t(1,y)\le\bar C t^{-3},\\
 \bar C=\frac{373901498071739108279096204237109367501674084730890663667211310806581045}{2^{256}}.
 \end{gathered}
 \tag{3.2}
\]
同时固定首坐标的Cauchy边缘和尾坐标符号对称性给
$\int_Fp_t\le t/\pi$。
因此
\[
 \norm{p_t(1,\cdot)}_{L^{5/4}([0,1]^4)}
 \le\bar C^{1/5}\pi^{-4/5}t^{1/5}.
\]
该面范数在$t=0$附近可积，Minkowski遂给小时间势在$F$上的$L^{5/4}$可积性。
这一估计包含全部$S_1,S_2$及交叠，
无需将四维的线外连续性结论移植到五维。

\subsection{小时间的正部尾比较}
记$R_\epsilon=G-H_\epsilon\ge0$、$D=\bar C/(2\epsilon^2)$、
$A=\epsilon^2/(2\pi)$。对任意$b\ge0$取
$t_0^2=\bar C/[2(b+D)]$，则(3.2)给
\[
 \int_F(R_\epsilon-b)_+\le\frac{AD}{b+D}.                 \tag{3.3}
\]
事实上$t_0$至$\epsilon$的点态积分至多$b$，
而零至$t_0$的面质量至多$t_0^2/(2\pi)$。
若$H_\epsilon\le M$于整个接收面$F$，则非减凸函数的正部积分公式给
\[
 \int_FG^{5/4}\le\int_FH_\epsilon^{5/4}+J_M,\qquad
 J_M=5A\int_0^1 z^2[D+(M-D)z^4]^{1/4}\dd z.             \tag{3.4}
\]
它来自将$(M+r)^{5/4}-M^{5/4}$代入(3.3)，再令
$z=(D/(D+b))^{1/4}$。需要的是非减凸性；任意凸函数不能直接使用此比较。
本证书取
\[
 M=\frac{488568596392024500738533860719}{250000000000000000000000000000}<3.
\]
该cap也有全闭面的解析证明。未分区峰值给所有$t>0$的
$p_t\le C_3t^{-3}+C_1t^{-1}$，五维Fourier峰值给
$p_t\le P_5t^{-5}$，其中$P_5=1/(65536\sqrt2\pi^3)$。
因此每个$T\ge\epsilon$都满足
\[
 \sup_F H_\epsilon\le\mathcal B(T):=
 \frac{C_3}{2}(\epsilon^{-2}-T^{-2})+
 C_1\log(T/\epsilon)+\frac{P_5}{4T^4}.
\]
$\mathcal B'(T)$的符号是$C_1T^4+C_3T^2-P_5$的符号，故仅有一个正stationary点。
原准确预算将系数换成有向端点，并包围该正根，得到上述有理$M$；
完整根及log余项在\path{upper_join132/}保留。
$M$仅为$F$上$H_{1/512}$的cap，不用于盒外锚点。
128个面板按$z^2$概率对凹四次根作Jensen上包；
全部257个根括界与128个面板和在第二轮被准确核验。
完整上包为$J_*$。所有无穷时间尾均保留在$H_\epsilon$中。

\section{空间下包：固定锚点与完整闭覆盖}
记$s_t(z)=t/[\pi(t^2+z^2)]$、$f_t(z)=\prod_{i=1}^{16}s_t(z_i)$。正纤维公式为
\[
 p_t(x)=\frac1{1024}\int_{\ker B}
 \prod_{j=1}^{16}\frac{t}{\pi[t^2+(B^Tx/16+w)_j^2]}\dd w.
\]
固定右逆和Cauchy的log-Hessian界，随后对固定锚点概率
$p_t(c)\dd t/H_\epsilon(c)$使用Jensen，得全空间位移公式
\[
 \log H_\epsilon(c+u)\ge\log H_\epsilon(c)
 +\nabla\log H_\epsilon(c)\cdot u-Q_\epsilon(c)|u|^2,\quad
 Q_\epsilon=\frac{\int_\epsilon^\infty p_t(c)t^{-2}\dd t}{16H_\epsilon(c)}.
 \tag{4.1}
\]
这不是沿空间路径冻结变化的后验。
正Gaussian混合给$t^5p_t(c)$非减；
与$t^{-2}$的反向Chebyshev比较在全尾及两段上分别应用，得到
\[
 Q_\epsilon\le\min\left\{q_{\rm old},\frac23A_\epsilon,
 A_\epsilon\left(\frac7{10}-\frac8{15}\frac{L}{U}\right)\right\},
 \quad A_\epsilon=\frac1{16\epsilon^2},                 \tag{4.2}
\]
其中$L\le H_{2\epsilon}(c)$、$H_\epsilon(c)\le U$。
$q_{\rm old}$只取同锚点、同截时、同物理度量的已认证上包。
比值前有负号，因而此处使用下端$L$和上端$U$。
没有$H_{2\epsilon}$证据时取$L=0$。

以$g_0=11/100000$为基准。
接收面内的位移满足$u_1=0$。每个锚点的准确$\Delta\le\log(H_\epsilon(c)/g_0)$、
四个物理坐标对数导数$\partial_j\log H_\epsilon(c)$的有向区间
$[\gamma_j^-,\gamma_j^+]$及$q\ge Q_\epsilon$给凹函数
\[
 \ell(u)=\Delta+\sum_{j=2}^5\min(\gamma_j^-u_j,\gamma_j^+u_j)-q|u|^2.
\]
同一锚点的全部16闭角点支付整个盒，包括锚点在盒外的情形。
$\epsilon=1/64$或$1/128$时均有$H_{1/128}\ge H_\epsilon$。
52360个有序根、58841次二分及111201个闭叶完整覆盖$F$。
全域最小余量$m>0$的完整分数和witness在准确证书中；
独立第二轮逐源、逐角点和二叉兄弟归并得到同一结果。

因$e^m\ge P_8(m)=\sum_{k=0}^8m^k/k!$，取
\[
 g_* =2^{-192}\left\lfloor2^{192}g_0P_8(m)\right\rfloor
 \le\essinf_FH_{1/128}.
 \tag{4.3}
\]
再用$G\ge H_{1/128}$、完整$G$的$-4$齐次性及方体对称性，才得到$G\ge g_*$于$Q$。
固定截时$H$没有被齐次外推。

\section{积分上包：几何正矩、时间弦与树支付}
\subsection{平移比与真实盒几何}
对面内位移$u_1=0$，$r=\norm u_\infty$，存在五方向提升$Ba=u$。
将尾坐标绝对值递减排列为$r=w_1\ge\cdots\ge w_4\ge w_5=0$，
沿相应符号链取
$a_0=-r/2$、$a_k=(w_k-w_{k+1})/2$。
单Cauchy密度的准确比值及asinh凹性给
\[
 \begin{gathered}
 \mathcal R(r/t)^{-1}\le p_t(c+u)/p_t(c)\le\mathcal R(r/t),\\
 \mathcal R(z)=
 (\sqrt{1+z^2/16}+z/4)^2(\sqrt{1+z^2/256}+z/16)^8.
 \end{gathered} \tag{5.1}
\]
并列、零坐标只产生零提升系数；$r=0$时比值为$1$。
该比较适用于全位移、全部正时间和所有锚点。

同一提升还满足$\sum|a_i|=r$、$|a_i|\le r/2$。
对无量纲score $S_i=t(\log s_t)'$和$T_i=t^2(\log s_t)''$，
当$r>0$时取$b_i=|a_i|/r$；$r=0$的二阶结论平凡。加权Cauchy不等式给
\[
 \frac{t^2D_a^2f_t}{r^2f_t}
 \le\sum_i b_i(S_i^2+b_iT_i).
\]
再对单Cauchy坐标的$q=(z_i/t)^2\ge0$使用
\[
 \frac{(4+2b)q-2b}{(1+q)^2}\le\frac{(b+2)^2}{4(b+1)}
 \le\frac{25}{24}\quad(0\le b\le1/2)
\]
给面内方向二阶上界
$D_u^2p_t\le(25/24)t^{-2}r^2p_t$。
这是二阶导数上界，不要求$p_t$凸。

令$U$在真实接收盒相对锚点的偏移上均匀分布，
$Z=\norm U_\infty$、$h=\max Z$、$\mu=\norm{\E U}_\infty$。
设$W(s)=\E(Z-s)_+$，
$J(z)=\int_0^hW(s)\mathcal R(sz)\dd s$。
Taylor公式、其非负余项上包及(5.1)得到
\[
 \frac{\E p_t(c+U)}{p_t(c)}
 \le\min\{\mathcal R(h/t),\,1+\mu/t+(25/24)t^{-2}J(1/t)\}. \tag{5.2}
\]
$J$非负、递增且凸，$J(0)=\E Z^2/2$。
若各偏移坐标为$[a_j,b_j]$，则
$\Pr(Z\le s)=\prod_j|[a_j,b_j]\cap[-s,s]|/(b_j-a_j)$，
故全部几何矩是分段多项式的准确有理积分。
偏心和盒外锚点的$\mu$必须保留。

\subsection{共同时间分布的加强}
在时间带$[a_t,b_t]$中令$a=1/a_t$、$b=1/b_t$，无穷尾取$b=0$。
由于$t^5p_t(c)$非减，(5.2)的递减时间函数可用概率
$4z^3\dd z/(a^4-b^4)$上包。记
\[
 E_j=\frac4{j+4}\frac{a^{j+4}-b^{j+4}}{a^4-b^4}.
\]
先对$J$取弦，再乘$z^2$，有
\[
 \E[z^2J(z)]\le\alpha J(a)+\beta J(b),\quad
 \alpha=\frac{E_3-bE_2}{a-b}>0,\quad
 \beta=\frac{aE_2-E_3}{a-b}>0.                         \tag{5.3}
\]
于是新均值因子为
$1+\mu E_1+(25/24)(\alpha J_{a,U}+\beta J_{b,U})$。
可先以$\min(J_{b,U},J_{a,U})$截紧右端。
$b=0$时$\alpha=4a^2/7$、$\beta=2a^2/21$，
使用准确$J(0)$，完整尾未截掉。
各条全区间有效上包分别积分后取min，再向上量化；不能对端点min作未经证明的插值。

\subsection{准确凸弦与完整森林}
由73个原始时间端点，带质量外端为
$H_L(a)-H_U(b)$与$H_U(a)-H_L(b)$；末端$H_\infty=0$。
正权相加给完整range $[L,U]$和均值上端。
对$\phi(x)=x^{5/4}$使用准确凸弦$\ell_{L,U}$，
\[
 \E\phi(H)\le\ell_{L,U}(\E H)\le\ell_{L,U}(m_U).
\]
这是全盒range上的弦支付，不是反向Jensen。
由$(x-1)^2(x^2+2x+3)\ge0$得准确弦斜率
\[
 s\ge U^{1/4}\left(1+\frac L{4U}\right).
\]
整数四次幂阈值给有理斜率下端$s_{\rm rat}$，不以舍入端点弦的斜率冒充它。
均值和积分均以$N=2^{160}$为单位；以下$m_{\rm old},m_{\rm new}$是均值上端的整数编码。
新的唯一总减量为
\[
 G_{\rm new}=\left\lfloor
 \frac{w(m_{\rm old}-m_{\rm new})s_{\rm rat}}{2^{4\ell}}\right\rfloor.
\]
它从原始已认证self扣除一次。
不能将两个相对同一基准的floor累加，也不能从已扣减self再扣总量。

有序根的轨道权为$w=4!/\prod m_j!$。
原树保持97972节点、52360根；叶支付self，非叶取
$\min(\mathrm{self},\sum\mathrm{childEff})$，仅对根求和。
真实canonical子盒及轨道权满足完整分割；
第二轮逐parent、incoming次数、可达性和16物理子格验证。
唯一未取新self的节点由其完整孩子支付。
最终准确根和$H_*$满足$\int_FH_{1/512}^{5/4}\le H_*$；
配合(3.4)得到$I\le H_*+J_*$。

\section{准确计算、两轮审计与复现}
原$H_\epsilon$和物理导数球来自已归档的完整有限扇：
280条生成射线、3504个五维锥。
每锥的Fourier积分为绝对行列式除五个resolvent因子的乘积；
在记录采用的公共归一化中，各时间系数为144。
令$D_c$为锚点$c$的正公共整数分母，$z_j$为满足$\psi(z_j)=144$的原整数射线，
$k_j=D_c\,c\cdot z_j$；每个锥使用其五个$k_j$。
具体地，以$Z_{\mathcal C}$表示各锥的五列射线矩阵，
\[
 H_\epsilon(c)=\frac1{(2\pi)^5}\operatorname{Re}
 \sum_{\mathcal C}|\det Z_{\mathcal C}|
 \int_\epsilon^\infty\prod_{j=1}^5(144t-i c\cdot z_j)^{-1}\dd t.
\]
这里$D_c$与小时间预算中的$D$不同。正截时积分以四阶差商计算，准确缩放为
\[ -\frac{D_c^4}{144(2\pi)^5}\,\mathrm{DD}_4\!\left[\frac12\log((144D_c\epsilon)^2+k^2)\right] \]
的加权和，其中每锥权重为$|\det Z_{\mathcal C}|$。
导数包含分母的商法则项及物理斜率$D_cz_{\rm tail}$。
原有向Arb/有理实现覆盖全部3504项，点的精度及generic守卫逐记录保存。
普通差商的锚点守卫没有删除接收域中的零或重复投影、边界及奇异点。

\begin{table}[ht]\centering\small
\begin{tabular}{p{.20\textwidth}p{.56\textwidth}r}
\toprule
实际批次&完整复核范围&秒\\ \midrule
166/169&97971组合总减量、97972节点、52360根；15个worker&2.55 / 2.60\\
第二轮178&97971$\times$73带、365因子、完整range/mean/幂/树及139全板&100.31\\
第二轮180&109446原记录、107530源、3812支撑、111201闭叶全角点&27.98\\
MATLAB R2026a&新完整常数的8项准确符号检查，私有批进程&96.90\\
\bottomrule
\end{tabular}
\caption{仅列本机实际运行证据。各数学批采用有限输入与可保存检查点，无程序墙钟上限。
秒数采用各程序自身登记，不代表累计CPU或全部研究耗时。
进程结束后已退出；没有关闭用户原有资源。}
\end{table}

首轮独立全链审计177及166增补通过后，第二轮178与180分别承担积分和全域空间复核，
主代理再逐项终审。178采用独立binary解析、CDF多项式、
512位根资格、Newton四次根和完整树；180采用exp-Taylor认证Delta、
不同分片及二叉兄弟归并。166、169、178、180的有限分片批各使用15个worker；
单标量检查与私有MATLAB另行登记。
没有将包完整性、哈希、操作记录或只准备未运行的checker当作新数学结果。

独立性有明确边界：原始$H_\epsilon$与四物理导数球仍是共享的、
已审查其模型及有向算术实现的输入；本轮没有重新运行第二套原始H引擎。
多个复核路径亦共享Python整数算术。
163、165、169的新组件具有原实际独审证据，第二轮审阅其理论、源码与接入，
没有伪称所有组件在每轮都重新生成。
完整资料、原始分数及来源SHA见仓库
\url{https://github.com/alizeisnotalice/Math-Research}的
\path{cube-dim-5-bounds/}，具体入口见附录。
旧三四维项目及原资料包仅作只读来源。

\section{结果范围}
定理\ref{thm:main}同时保留Lebesgue背景、中心轴平行方体、共同半径、
一般$L^1$及全部正尺度。它严格改进已核查公开、原内部及接续的上界32。
完整新旧比较使用准确四次幂与有理根区间；较小模型成本只有在完整
$g,I$与分析传递均闭合后才转化为算子界。
局部$g_{115}$候选、后续band曲率与更细斜率准备未进入本定理。
本稿没有提高冻结接续下界，也没有据此触发原双侧目标的完成声明。

\appendix
\section{短整数证书与完整准确定义}
短证书取
\[
 N=10^{20},\quad G_0=11231606208775453,\quad
 I_0=4474860503582229,\quad K=28936079.
\]
从完整准确分数向外取整得到$g_L=G_0/N\le g_*$、
$I\le I_U=I_0/N$。纯整数严格门为
\[
 K^4G_0^5-3125I_0^4\,10^{44}>0.
\]
故命题\ref{prop:transfer}直接给$C_5<K/10^6$。
该门由独立整数解析和交叉乘法再次实际核验；没有从小数构数。
读者可作最小标量复核：
\begin{lstlisting}[language=Python]
G = 11231606208775453
I = 4474860503582229
K = 28936079
assert K**4 * G**5 > 3125 * I**4 * 10**44
\end{lstlisting}
此四行仅复核最后的标量门；完整$g,I$资格由正文及数据证书支付。

完整定义为
\[
 \begin{gathered}
 g_*=\frac{705019348242842195339748239684514494542987839533613919}{2^{192}},
 \\
 H_*=\frac{3994994121862399564161069833676547838746859}
 {91343852333181432387730302044767688728495783936}.
 \end{gathered}
\]
\[
\resizebox{.98\textwidth}{!}{$
 J_*=
 \frac{17096357665258533257620191768869876487723650524339571180562654263529837424217770456490029292234384306465669120}
 {16879779982267146520848231959882343332105778642507237944706013844469724284863538912850515249620229857207350532657943}.
$}
\]
原接纳证书为\path{round019/primary_upper166_acceptance.json}；
其新旧四次幂差、根区间、三类32门槛比较以及完整来源SHA全部随文件给出。

\section{复核入口与证据分类}
\begin{longtable}{p{.35\textwidth}p{.55\textwidth}}
\toprule 位置&用途及边界\\ \midrule\endhead
\path{primary/finite_source_transfer.md}&有限源、同一接收集、一般输入和全尺度证明。\\
\path{full_audit177/}&首轮完整审计；\path{addendum166/}补齐组合积分资格。\\
\path{second_audit178/}&第二轮全量积分reader、真实结果及短常数独审加页。\\
\path{strict_g180/portable/}&第二轮全域空间reader及仓库数据参数视图；
原始raw操作回连继承本轮实际31006，具体说明见其README。\\
\path{upper_integral167/packet/}&完整原H端点输入。传输认证与数学认证严格区分。\\
\path{upper_combined_archive179/}&组合积分数学JSON的去重字节恢复、来源映射与完整hash。\\
\path{results/bounds_ledger.json}&完整准确账本、公开与内部比较、历史界及当前状态。\\
\path{round019/final_upper_review.json}&主终审、共享依赖边界及本稿生成条件。\\
\bottomrule
\end{longtable}

数值候选、局部条件、完整严格证书、准确排除与计算故障在账本中分别记录。
本稿不把准备期路径或schema失败当作数学反例；
本轮两项第二轮全量reader均真实完成。
为防止覆盖冻结证据，复核器使用新的输出目录；原始过程的PID和调用记录
只证明其历史运行，不替代本次实际结果。

\begin{thebibliography}{9}
\bibitem{Aldaz2000}
J. M. Aldaz. A remark on the centered $n$-dimensional Hardy--Littlewood maximal function.
\emph{Czechoslovak Mathematical Journal}, 50 (2000), 103--112.
\bibitem{Aldaz2011}
J. M. Aldaz. The weak type $(1,1)$ bounds for the maximal function associated to cubes
grow to infinity with the dimension.
\emph{Annals of Mathematics}, 173 (2011), 1013--1023.
\href{https://doi.org/10.4007/annals.2011.173.2.10}{doi:10.4007/annals.2011.173.2.10}.
\bibitem{Aldaz2020}
J. M. Aldaz. Kissing numbers and the centered maximal operator.
arXiv:1811.10478v2, 2020.
\url{https://arxiv.org/abs/1811.10478}.
\bibitem{Lin2026}
Y. Lin. Two partial balayage principles for weak-type estimates.
作者原始方法网页，2026年10月25日至26日核查。
\url{https://coolrmal.github.io/articles/two-partial-balayage-principles/}.
\end{thebibliography}
\end{document}
